\documentclass[reqno]{amsart}
\usepackage[T1]{fontenc}
\usepackage{amsmath,amssymb,amsthm}
\usepackage{xr-hyper}
\usepackage[hidelinks]{hyperref}
\externaldocument[sm-]{20261125T000000--slope-method__workbook}
\externaldocument[dt-]{20261125T002000--definite-theta-lifts__workbook}
\externaldocument[cp-]{20261125T006500--covolumes-and-periods__workbook}
\externaldocument[as-]{20261125T021500--assembly-at-one-scale__workbook}

% Chapter number; \ref{chap} gives it.
\newcounter{chap}
% Chapter numbers, generated from 20261125T012000--grh-workbook__guide.org
% by scripts/make-chapters.py.  Do not edit: change the order of the
% index and rerun the script (./preview-site.sh and CI do this).
%
% Usage in a sheet, after \newcounter{chap}:
%   % Chapter numbers, generated from 20261125T012000--grh-workbook__guide.org
% by scripts/make-chapters.py.  Do not edit: change the order of the
% index and rerun the script (./preview-site.sh and CI do this).
%
% Usage in a sheet, after \newcounter{chap}:
%   % Chapter numbers, generated from 20261125T012000--grh-workbook__guide.org
% by scripts/make-chapters.py.  Do not edit: change the order of the
% index and rerun the script (./preview-site.sh and CI do this).
%
% Usage in a sheet, after \newcounter{chap}:
%   \input{chapters}\chapnumber{<sheet ID>}
\makeatletter
\providecommand\chapters@set[2]{\@namedef{chapters@#1}{#2}}
\chapters@set{20261125T001000}{1}% zeros-to-ramp-sums
\chapters@set{20261125T001100}{2}% inert-primes-and-estimate-b
\chapters@set{20261125T005000}{3}% galois-fourier-and-averaging
\chapters@set{20261125T005500}{4}% columns-from-a-sphere
\chapters@set{20261125T000000}{5}% slope-method
\chapters@set{20261126T030100}{6}% a-twist-for-each-prime
\chapters@set{20261125T003000}{7}% the-chain-and-the-square-root
\chapters@set{20261125T010000}{8}% rank-lemma-and-zero-estimate
\chapters@set{20261125T010500}{9}% poincare-series-and-near-orthogonality
\chapters@set{20261124T223000}{10}% elliptic-curves-for-the-proof
\chapters@set{20261125T009500}{11}% fock-model-and-theta-intertwiners
\chapters@set{20261124T223100}{12}% complex-multiplication-and-inert-reduction
\chapters@set{20261125T008000}{13}% one-cm-determinant
\chapters@set{20261124T214500}{14}% supersingular-theta-divisibility
\chapters@set{20261124T214600}{15}% the-isotropic-plane
\chapters@set{20261125T013700}{16}% pairings-on-a-formal-group-of-height-two
\chapters@set{20261125T013500}{17}% hecke-nondegeneracy-on-the-tree
\chapters@set{20261126T030200}{18}% the-bound-is-attained
\chapters@set{20261125T009000}{19}% kronecker-limit-and-chowla-selberg
\chapters@set{20261125T002000}{20}% definite-theta-lifts
\chapters@set{20261125T015500}{21}% siegel-weil-in-the-smallest-case
\chapters@set{20261125T002100}{22}% definite-theta-lifts-and-norms
\chapters@set{20261125T016500}{23}% hecke-operators-and-theta-lifts
\chapters@set{20261125T016700}{24}% divisibility-along-the-chain
\chapters@set{20261125T004000}{25}% denominators-and-sha
\chapters@set{20261125T004100}{26}% tate-twists-and-sha
\chapters@set{20261126T020430}{27}% the-minus-part-by-gauss-sums
\chapters@set{20261125T014100}{28}% congruence-modules
\chapters@set{20261125T014500}{29}% eisenstein-congruences-and-selmer-classes
\chapters@set{20261125T014700}{30}% pseudocharacters-and-the-length-bound
\chapters@set{20261125T011000}{31}% galois-representations-black-box
\chapters@set{20261126T030300}{32}% the-data
\chapters@set{20261126T020210}{33}% bessel-poincare-series-and-their-gram-operators
\chapters@set{20261126T030400}{34}% the-seesaw-identity
\chapters@set{20261126T030500}{35}% divisibility-of-the-composite-lift
\chapters@set{20261126T030600}{36}% theorem-a-short-forms
\chapters@set{20261126T030700}{37}% the-determinant-of-a-chain
\chapters@set{20261126T030800}{38}% the-cost-of-one-lift
\chapters@set{20261126T030900}{39}% denominators-of-projections-to-lifts
\chapters@set{20261126T031000}{40}% twists-with-few-congruences
\chapters@set{20261126T031100}{41}% theorem-c-and-the-main-inequality
\providecommand\chapnumber[1]{%
  \@ifundefined{chapters@#1}%
    {\PackageError{chapters}{Sheet #1 is not listed in the index}%
      {List it in 20261125T012000--grh-workbook__guide.org\MessageBreak
       and run scripts/make-chapters.py.}}%
    {\edef\thechap{\@nameuse{chapters@#1}}}}
\makeatother
\chapnumber{<sheet ID>}
\makeatletter
\providecommand\chapters@set[2]{\@namedef{chapters@#1}{#2}}
\chapters@set{20261125T001000}{1}% zeros-to-ramp-sums
\chapters@set{20261125T001100}{2}% inert-primes-and-estimate-b
\chapters@set{20261125T005000}{3}% galois-fourier-and-averaging
\chapters@set{20261125T005500}{4}% columns-from-a-sphere
\chapters@set{20261125T000000}{5}% slope-method
\chapters@set{20261126T030100}{6}% a-twist-for-each-prime
\chapters@set{20261125T003000}{7}% the-chain-and-the-square-root
\chapters@set{20261125T010000}{8}% rank-lemma-and-zero-estimate
\chapters@set{20261125T010500}{9}% poincare-series-and-near-orthogonality
\chapters@set{20261124T223000}{10}% elliptic-curves-for-the-proof
\chapters@set{20261125T009500}{11}% fock-model-and-theta-intertwiners
\chapters@set{20261124T223100}{12}% complex-multiplication-and-inert-reduction
\chapters@set{20261125T008000}{13}% one-cm-determinant
\chapters@set{20261124T214500}{14}% supersingular-theta-divisibility
\chapters@set{20261124T214600}{15}% the-isotropic-plane
\chapters@set{20261125T013700}{16}% pairings-on-a-formal-group-of-height-two
\chapters@set{20261125T013500}{17}% hecke-nondegeneracy-on-the-tree
\chapters@set{20261126T030200}{18}% the-bound-is-attained
\chapters@set{20261125T009000}{19}% kronecker-limit-and-chowla-selberg
\chapters@set{20261125T002000}{20}% definite-theta-lifts
\chapters@set{20261125T015500}{21}% siegel-weil-in-the-smallest-case
\chapters@set{20261125T002100}{22}% definite-theta-lifts-and-norms
\chapters@set{20261125T016500}{23}% hecke-operators-and-theta-lifts
\chapters@set{20261125T016700}{24}% divisibility-along-the-chain
\chapters@set{20261125T004000}{25}% denominators-and-sha
\chapters@set{20261125T004100}{26}% tate-twists-and-sha
\chapters@set{20261126T020430}{27}% the-minus-part-by-gauss-sums
\chapters@set{20261125T014100}{28}% congruence-modules
\chapters@set{20261125T014500}{29}% eisenstein-congruences-and-selmer-classes
\chapters@set{20261125T014700}{30}% pseudocharacters-and-the-length-bound
\chapters@set{20261125T011000}{31}% galois-representations-black-box
\chapters@set{20261126T030300}{32}% the-data
\chapters@set{20261126T020210}{33}% bessel-poincare-series-and-their-gram-operators
\chapters@set{20261126T030400}{34}% the-seesaw-identity
\chapters@set{20261126T030500}{35}% divisibility-of-the-composite-lift
\chapters@set{20261126T030600}{36}% theorem-a-short-forms
\chapters@set{20261126T030700}{37}% the-determinant-of-a-chain
\chapters@set{20261126T030800}{38}% the-cost-of-one-lift
\chapters@set{20261126T030900}{39}% denominators-of-projections-to-lifts
\chapters@set{20261126T031000}{40}% twists-with-few-congruences
\chapters@set{20261126T031100}{41}% theorem-c-and-the-main-inequality
\providecommand\chapnumber[1]{%
  \@ifundefined{chapters@#1}%
    {\PackageError{chapters}{Sheet #1 is not listed in the index}%
      {List it in 20261125T012000--grh-workbook__guide.org\MessageBreak
       and run scripts/make-chapters.py.}}%
    {\edef\thechap{\@nameuse{chapters@#1}}}}
\makeatother
\chapnumber{<sheet ID>}
\makeatletter
\providecommand\chapters@set[2]{\@namedef{chapters@#1}{#2}}
\chapters@set{20261125T001000}{1}% zeros-to-ramp-sums
\chapters@set{20261125T001100}{2}% inert-primes-and-estimate-b
\chapters@set{20261125T005000}{3}% galois-fourier-and-averaging
\chapters@set{20261125T005500}{4}% columns-from-a-sphere
\chapters@set{20261125T000000}{5}% slope-method
\chapters@set{20261126T030100}{6}% a-twist-for-each-prime
\chapters@set{20261125T003000}{7}% the-chain-and-the-square-root
\chapters@set{20261125T010000}{8}% rank-lemma-and-zero-estimate
\chapters@set{20261125T010500}{9}% poincare-series-and-near-orthogonality
\chapters@set{20261124T223000}{10}% elliptic-curves-for-the-proof
\chapters@set{20261125T009500}{11}% fock-model-and-theta-intertwiners
\chapters@set{20261124T223100}{12}% complex-multiplication-and-inert-reduction
\chapters@set{20261125T008000}{13}% one-cm-determinant
\chapters@set{20261124T214500}{14}% supersingular-theta-divisibility
\chapters@set{20261124T214600}{15}% the-isotropic-plane
\chapters@set{20261125T013700}{16}% pairings-on-a-formal-group-of-height-two
\chapters@set{20261125T013500}{17}% hecke-nondegeneracy-on-the-tree
\chapters@set{20261126T030200}{18}% the-bound-is-attained
\chapters@set{20261125T009000}{19}% kronecker-limit-and-chowla-selberg
\chapters@set{20261125T002000}{20}% definite-theta-lifts
\chapters@set{20261125T015500}{21}% siegel-weil-in-the-smallest-case
\chapters@set{20261125T002100}{22}% definite-theta-lifts-and-norms
\chapters@set{20261125T016500}{23}% hecke-operators-and-theta-lifts
\chapters@set{20261125T016700}{24}% divisibility-along-the-chain
\chapters@set{20261125T004000}{25}% denominators-and-sha
\chapters@set{20261125T004100}{26}% tate-twists-and-sha
\chapters@set{20261126T020430}{27}% the-minus-part-by-gauss-sums
\chapters@set{20261125T014100}{28}% congruence-modules
\chapters@set{20261125T014500}{29}% eisenstein-congruences-and-selmer-classes
\chapters@set{20261125T014700}{30}% pseudocharacters-and-the-length-bound
\chapters@set{20261125T011000}{31}% galois-representations-black-box
\chapters@set{20261126T030300}{32}% the-data
\chapters@set{20261126T020210}{33}% bessel-poincare-series-and-their-gram-operators
\chapters@set{20261126T030400}{34}% the-seesaw-identity
\chapters@set{20261126T030500}{35}% divisibility-of-the-composite-lift
\chapters@set{20261126T030600}{36}% theorem-a-short-forms
\chapters@set{20261126T030700}{37}% the-determinant-of-a-chain
\chapters@set{20261126T030800}{38}% the-cost-of-one-lift
\chapters@set{20261126T030900}{39}% denominators-of-projections-to-lifts
\chapters@set{20261126T031000}{40}% twists-with-few-congruences
\chapters@set{20261126T031100}{41}% theorem-c-and-the-main-inequality
\providecommand\chapnumber[1]{%
  \@ifundefined{chapters@#1}%
    {\PackageError{chapters}{Sheet #1 is not listed in the index}%
      {List it in 20261125T012000--grh-workbook__guide.org\MessageBreak
       and run scripts/make-chapters.py.}}%
    {\edef\thechap{\@nameuse{chapters@#1}}}}
\makeatother
\chapnumber{20261125T006000}
\renewcommand{\thesection}{\thechap.\arabic{section}}

\theoremstyle{definition}
\newtheorem{problem}{Problem}[section]
\theoremstyle{plain}
\newtheorem{theorem}{Theorem}[section]

\newcommand{\Z}{\mathbb{Z}}
\newcommand{\Q}{\mathbb{Q}}
\newcommand{\R}{\mathbb{R}}
\newcommand{\C}{\mathbb{C}}
\newcommand{\Pb}{\mathbb{P}}
\newcommand{\Oc}{\mathcal{O}}
\newcommand{\Hc}{\mathcal{H}}
\newcommand{\Ac}{\mathcal{A}}
\newcommand{\ii}{\mathrm{i}}
\newcommand{\mumax}{\mu_{\max}}
\newcommand{\mumin}{\mu_{\min}}
\newcommand{\covol}{\operatorname{covol}}
\newcommand{\Cl}{\mathrm{Cl}}
\newcommand{\Sym}{\operatorname{Sym}}
\newcommand{\hF}{h_{\mathrm{F}}}
\newcommand{\cc}{\hat{c}_1}

\setlength{\emergencystretch}{2em}

\title{Covolumes of lattices of sections}
\author{}
\date{November 25, 2026}

\begin{document}
\maketitle
\refstepcounter{chap}\label{chap}

\paragraph{Prerequisites.}
Euler--Maclaurin summation and Stirling's formula.  Arithmetic intersection theory is used only through the statement of Theorem~\ref{thm:HS}.  From Chapter~\ref{sm-chap}, \emph{The slope method}: Section~\ref{sm-sec:slopes} (normed spaces, degrees and slopes, \(\mu^\sharp\)), and the definition of \(\mumax\) in Section~\ref{sm-sec:image}.  From Chapter~\ref{dt-chap}, \emph{Masses and the count of forms}: Sections~\ref{dt-sec:lattices}--\ref{dt-sec:forms}, in Section~\ref{sec:definite}.

The slope method of Chapter~\ref{sm-chap} needs a lower bound for the slope of its source, a lattice of sections or of automorphic forms.  In the argument the lattice of automorphic forms is the one whose dual is the target of the maps \cite[\S15.4]{Adr}.  How is its covolume computed?  We compute it exactly for polynomials on a projective space and for forms of weight \((a)\) on a definite unitary group of rank at most two, compare with the arithmetic Hilbert--Samuel theorem, and ask which uniformity the argument needs.  The period that enters through the integral structure is the subject of Chapter~\ref{cp-chap}.  Degrees and slopes are those of Section~\ref{sm-sec:slopes}; over \(\Q\), \(\deg L = -\log\covol(L)\) (Problem~\ref{sm-pr:deg-Q}).

\section{Projective spaces}\label{sec:P1}

Let \(E_D \subset \Z[x, y]\) be the lattice of binary forms of degree \(D\) with the Fubini--Study \(L^2\) norm of Section~\ref{sm-sec:slopes}.  The monomials are orthogonal, with \(\lVert x^a y^{D-a} \rVert^2 = a!\,(D-a)!/(D+1)!\) (Problem~\ref{sm-pr:FS}), so the covolume has a closed form.

\begin{problem}\label{pr:degED}
\(\displaystyle \deg E_D = \tfrac12 (D+1)\log(D+1) + \tfrac12 \sum_{k=1}^{D} (2k - D - 1) \log k .\)
\end{problem}

\begin{problem}\label{pr:degED-asym}
\(\deg E_D = \tfrac14 D^2 + \tfrac14 D\log D + \tfrac12\bigl(1 - \tfrac12\log 2\pi\bigr) D + O(\log D)\).  For \(D = 10\) the exact value is \(\deg E_{10} = 32.1969\ldots\)
\end{problem}

Problems~\ref{pr:degED} and~\ref{pr:degED-asym} prove the following, which refines Problem~\ref{sm-pr:HS-P1}.

\begin{theorem}\label{thm:P1}
The lattice \(H^0(\Pb^1_\Z, \Oc(D)) = E_D\), with the \(L^2\) norm of the Fubini--Study metric and measure, has
\[
  \deg E_D = \tfrac14 D^2 + \tfrac14 D\log D + O(D) .
\]
\end{theorem}

The term \(\tfrac14 D\log D\) is the sum of \(\tfrac12 D\log D\) from the factor \(D+1\) in \(\lVert x^a y^{D-a} \rVert^{-2} = (D+1)\binom{D}{a}\), which comes from the normalization of the measure, and \(-\tfrac14 D\log D\) from the binomial coefficients.  How far do the slopes of sublattices spread?  Recall that \(\mumax(L)\) is the largest slope of a nonzero sublattice of \(L\), and let \(\mumin(L)\) be the smallest slope of a nonzero quotient of \(L\).

\begin{problem}\label{pr:orth-sum}
Let \(L = L_1 \oplus L_2\) be an orthogonal sum of Euclidean lattices.  Then \(\mumax(L) = \max(\mumax(L_1), \mumax(L_2))\).
\end{problem}

\begin{problem}\label{pr:hn-P1}
\(\mumax(E_D) = \tfrac12 \log\bigl((D+1)\binom{D}{\lfloor D/2 \rfloor}\bigr) = \tfrac12 D\log 2 + O(\log D)\), and \(\mumin(E_D) = \tfrac12\log(D+1)\).
\end{problem}

A lower bound for \(\mu(E_D)\) is all the slope method uses, so the sublattices of large slope cost nothing.  Let \(\lVert f \rVert_\infty := \sup_z \lvert f(z, 1) \rvert (1 + \lvert z \rvert^2)^{-D/2}\).

\begin{problem}\label{pr:bergman}
For \(f \in E_D \otimes \C\), \(\lVert f \rVert \le \lVert f \rVert_\infty \le (D+1)^{1/2} \lVert f \rVert\).  The second inequality is an equality for \(f = (\bar w x + y)^D\).
\end{problem}

In \(n\) variables, let \(E_{n,a} \subset \Z[x_1, \dots, x_n]\) be the lattice of forms of degree \(a\), with \(\lVert f \rVert^2 := \int_{S^{2n-1}} \lvert f \rvert^2 \, d\sigma\) for the invariant probability measure \(\sigma\) on the unit sphere of \(\C^n\).  For \(n = 2\) this is the norm of \(E_a\); in general it is the Fubini--Study \(L^2\) norm on \(\Pb^{n-1}\).

\begin{problem}\label{pr:PN}
The monomials are orthogonal in \(E_{n,a}\), and \(\lVert x^\nu \rVert^2 = \nu!\,(n-1)!/(a+n-1)!\).
\end{problem}

\begin{problem}\label{pr:base-change}
Let \(V\) be a normed space over \(K\), and \(K'/K\) a finite extension.  Give \(V \otimes_K K'\) the unit lattices \(L_v \otimes \Oc_{K'_w}\) and the norms extended to \(V \otimes_K K'_w\).  Then \(\mu^\sharp(V \otimes_K K') = \mu^\sharp(V)\).  In particular \(\mu^\sharp(E_D \otimes \Oc_K) = 2\mu(E_D)\).
\end{problem}

\section{The arithmetic Hilbert--Samuel theorem}\label{sec:HS}

Let \(X\) be a projective flat integral scheme over \(\Z\) of dimension \(d+1\) with smooth generic fiber, \(\bar L\) an ample line bundle on \(X\) with a smooth, conjugation-invariant Hermitian metric of positive curvature, and \(dV\) a smooth conjugation-invariant volume form on \(X(\C)\).  Give \(H^0(X, L^{\otimes n})\) the norm \(\lVert s \rVert^2 := \int_{X(\C)} \lvert s \rvert^2 \, dV\), and let \(\cc(\bar L)^{d+1}\) be the arithmetic self-intersection of Gillet and Soul\'e \cite{GS90}.

\begin{theorem}[Gillet--Soul\'e, Abbes--Bouche]\label{thm:HS}
As \(n \to \infty\),
\[
  \deg H^0(X, L^{\otimes n}) = \frac{\cc(\bar L)^{d+1}}{(d+1)!}\, n^{d+1} + o(n^{d+1}) .
\]
\end{theorem}

This is due to Gillet and Soul\'e for regular \(X\) \cite{GS92} and to Abbes and Bouche in general \cite[\S1, Th\'eor\`eme principal]{AB95}; the error depends on \(X\), \(\bar L\) and \(dV\), and is not explicit.  For \(\Pb^N_\Z\) with the Fubini--Study metric, Bost, Gillet and Soul\'e computed
\begin{equation}\label{eq:hPN}
  \cc(\bar\Oc(1))^{N+1} = \frac12 \sum_{k=1}^{N} \sum_{m=1}^{k} \frac1m
\end{equation}
\cite{BGS}.

\begin{problem}\label{pr:HS-P1-check}
Theorem~\ref{thm:P1} and Theorem~\ref{thm:HS} give \(\cc(\bar\Oc(1))^2 = \tfrac12\) on \(\Pb^1_\Z\), in agreement with~\eqref{eq:hPN}.  Replacing the metric of \(\Oc(1)\) by \(e^{-c}\lvert\cdot\rvert\) adds \(cD(D+1)\) to \(\deg E_D\) and \(2c\) to \(\cc(\bar\Oc(1))^2\).
\end{problem}

\begin{problem}\label{pr:L2-sup}
Let \(\lVert\cdot\rVert_1\) and \(\lVert\cdot\rVert_2\) be Euclidean norms on a lattice \(L\) of rank \(r\), with \(\lVert v \rVert_1 \le C\lVert v \rVert_2\) for all \(v\).  Then \(\deg_1 L \ge \deg_2 L - r\log C\).  Hence, since \(\operatorname{rank} H^0(X, L^{\otimes n}) = O(n^d)\), the leading term in Theorem~\ref{thm:HS} does not depend on \(dV\).
\end{problem}

\section{Definite unitary groups}\label{sec:definite}

On a definite group the Shimura variety is a finite set, and the covolume of the lattice of integral forms can be computed by hand.  The argument needs the covolume on unitary groups of signature \((p, q)\) with \(p, q > 0\), whose forms have discrete series at infinity and whose Shimura varieties have positive dimension \cite[\S15.14 and Hyp.~15.58]{Adr}; the definite case is a model.  We use the notation of Chapter~\ref{dt-chap}, \emph{Masses and the count of forms}: \(\Oc = \Z[\ii]\), the class set \(\Cl_n\), the groups \(\Gamma_L\), the mass \(M_n\), the spaces \(V_\lambda(L)\), and the forms \(\Ac_\lambda(n)\) with the norm \(\lVert\varphi\rVert^2 = \sum_L \lvert\Gamma_L\rvert^{-1}\lVert\varphi(L)\rVert^2\).  Let \(V_\lambda(L)_\Z\) be the image of \(L^{\otimes\lvert\lambda\rvert}\) under the orthogonal projection onto \(V_\lambda(L)\); for \(\lambda = (a)\) it is the \(\Oc\)-span of the \(e^\nu\).  The \emph{integral forms} make up the \(\Oc\)-lattice \(\Ac_\lambda(n)_\Z\) of \(\varphi\) with \(\varphi(L) \in V_\lambda(L)_\Z\) for every \(L\).  Degrees are over \(\Q(\ii)\), so \(\mu^\sharp = \mu\).

\begin{problem}\label{pr:def-deg}
\(\displaystyle \deg \Ac_\lambda(n)_\Z = \sum_{L \in \Cl_n} \Bigl( \deg V_\lambda(L)_\Z^{\Gamma_L} + r_L \log\lvert\Gamma_L\rvert \Bigr)\), where \(r_L\) is the rank of the invariants \(V_\lambda(L)_\Z^{\Gamma_L}\).
\end{problem}

\begin{problem}\label{pr:def-n1}
For \(n = 1\) and \(4 \mid a\), \(\Ac_{(a)}(1)_\Z\) has rank one and slope \(\log 4 = -\log M_1\).
\end{problem}

\begin{problem}\label{pr:sym-sphere}
With the norm of Problem~\ref{dt-pr:sym-norm}, \(\lVert e^\nu \rVert^2 = \binom{a+n-1}{n-1} \lVert x^\nu \rVert^2\), where \(\lVert x^\nu \rVert\) is the norm of Problem~\ref{pr:PN}.  Hence \(V_{(a)}(I_n)_\Z = E_{n,a} \otimes \Oc\), with its norm multiplied by \(\binom{a+n-1}{n-1}^{1/2}\).
\end{problem}

For \(n = 2\) the class set is \(\{I_2\}\), and \(\lvert\Gamma_{I_2}\rvert = 32\) (Problems~\ref{dt-pr:class-one} and~\ref{dt-pr:aut}).

\begin{problem}\label{pr:def-n2}
Let \(4 \mid a\).  The invariants \(V_{(a)}(I_2)_\Z^{\Gamma}\) have the basis \(\sum_{\mu \in \{\nu, \nu'\}} e^\mu\), where \(\nu = (\nu_1, \nu_2)\) runs over the pairs with \(4 \mid \nu_1\) and \(\nu_1 \le \nu_2\), and \(\nu' = (\nu_2, \nu_1)\).  Consequently
\[
  \deg \Ac_{(a)}(2)_\Z = \sum_{\nu} \Bigl( \log 32 - \log\lvert\{\nu, \nu'\}\rvert + \log\tbinom{a}{\nu_1} \Bigr),
\]
and \(\deg\Ac_{(8)}(2)_\Z = \log 35840\).
\end{problem}

\begin{problem}\label{pr:def-n2-asym}
\(\mu\bigl(\Ac_{(a)}(2)_\Z\bigr) = \tfrac12 a + O(\log a)\) as \(a \to \infty\) through multiples of \(4\).
\end{problem}

So for \(n = 2\) and \(\lambda = (a)\) the rank is about \(4M_2\dim V_{(a)}\) (Problem~\ref{dt-pr:forms-2}), each form contributes \(\log 32 = -\log M_2\) through the stabilizer, and the slope is that of \(E_a \otimes \Oc\).  We do not treat larger \(n\).

\section{Uniformity}\label{sec:uniform}

Theorem~\ref{thm:HS} concerns a fixed variety and a growing power of a fixed line bundle.  In Alpöge's argument \cite{Alp}, \cite{Zh} the varieties are Shimura varieties of unitary groups, or flag bundles over them, whose dimension grows; the level and the field vary; and the weight is given rather than tending to infinity.  By Problem~\ref{pr:base-change} a bound for \(\mu^\sharp\) can be uniform in the field.  The other parameters are harder.

\begin{problem}\label{pr:rank-ratio}
\(\displaystyle 1 \le \binom{a+n-1}{n-1} \Bigm/ \frac{a^{n-1}}{(n-1)!} \le e^{n^2/2a}\), and the ratio tends to infinity if \(n^2/a\) does.
\end{problem}

So the leading term \(a^{n-1}/(n-1)!\) of the Hilbert polynomial approximates the rank \(\binom{a+n-1}{n-1}\) of \(E_{n,a}\) only when \(a\) is large compared with \(n^2\), the square of the dimension.

\begin{problem}\label{pr:trivial-lower}
\(\mu(E_{n,a}) \ge \tfrac12\log\binom{a+n-1}{n-1}\).
\end{problem}

This bound is uniform but misses the main term: for fixed \(n\), \eqref{eq:hPN} gives \(\mu(E_{n,a}) \sim c_n a\) with \(c_n = \tfrac{1}{2n}\sum_{k<n}\sum_{m \le k} 1/m\).  Chow forms give lower bounds in which every constant is explicit.  For an irreducible \(X \subset \Pb^n_\Q\) of dimension \(r - 1\) with homogeneous ideal \(I \subset \Z[x_0, \dots, x_n]\), David and Philippon bound \(-\log\covol\bigl((\Z[x_0, \dots, x_n]/I)_a\bigr)\) from below by \(h(I)\,\bigl((a - c)/r\bigr)^r\) minus an explicit error, for every \(a\) above an explicit threshold \(c\), where \(h(I)\) is a height of \(X\) \cite[Prop.~4.10 and Cor.~4.11]{DPh}.  They do not aim at the leading term: the factor \(r^{-r}\) stands where Theorem~\ref{thm:HS} has \(1/r!\), which is larger by a factor exponential in \(r\).

\begin{problem}\label{pr:level}
Let \(N \ge 1\), and let \(E_D(N) \subseteq E_D\) be the sublattice of \(f\) whose coefficient of \(x^D\) is divisible by \(N\), that is, \(f(1,0) \equiv 0 \pmod N\).  Then \(\mu(E_D(N)) = \mu(E_D) - (\log N)/(D+1)\).
\end{problem}

For automorphic forms a level structure also removes the terms \(r_L\log\lvert\Gamma_L\rvert\) of Problem~\ref{pr:def-deg} and multiplies the rank by the index.

\paragraph{Questions.}
In \cite{Adr} the covolume bound the argument needs is \cite[Hyp.~15.58]{Adr}; see the ledger in Section~\ref{as-sec:map} of Chapter~\ref{as-chap}.  \cite[\S4.3]{Zh} describes Alpöge's argument as applying Theorem~\ref{thm:HS} while keeping track of the losses from changing the weight and the level and from the boundary of the compactification.  The two sources apply it to different lattices: \cite[\S4.3]{Zh} to the initial lattices of the construction, the base lattices of Chapter~\ref{as-chap}, and \cite[\S15]{Adr}, which the workbook follows, to the target, with the base lattices treated by a computation with Weyl modules; see the remark after the ledger in Section~\ref{as-sec:map} of Chapter~\ref{as-chap}.  We do not know a published form of Theorem~\ref{thm:HS} that is uniform in all the parameters of this section.  We also do not know whether Problems~\ref{pr:def-n2} and~\ref{pr:def-n2-asym} extend to \(\Ac_\lambda(n)_\Z\) for all \(n\) with an error \(O(\log\lvert\lambda\rvert)\) per form, uniformly in \(n\).

\section*{Hints}
\paragraph{Section~\ref{sec:P1}.}
\begin{description}
\item[\ref{pr:degED}] Problem~\ref{sm-pr:FS}; collect the terms \(\log k\).
\item[\ref{pr:degED-asym}] Expand \(\sum_{k \le D} k\log k\) and \(\log D!\) by Euler--Maclaurin, to \(O(\log D)\).  Only Stirling's constant \(\tfrac12\log 2\pi\) survives, multiplied by \(-(D+1)/2\).
\item[\ref{pr:orth-sum}] For a sublattice \(V \subseteq L\), project to \(L_2\) along \(L_1\).  The quotient norm on \(V/(V \cap L_1)\) is at least the norm of the projection.  Use Problem~\ref{sm-pr:additivity}.
\item[\ref{pr:hn-P1}] Apply Problem~\ref{pr:orth-sum} to the monomial lines.  For \(\mumin\), pass to the dual lattice.
\item[\ref{pr:bergman}] Expand \(f\) in the orthonormal basis of normalized monomials and apply Cauchy--Schwarz at the point \(z\).  The sum of the squares of the basis functions at \(z\) is \(D+1\) times \((1 + \lvert z \rvert^2)^D\).
\item[\ref{pr:PN}] Integrate \(\lvert z^\nu \rvert^2 e^{-\lvert z \rvert^2}\) over \(\C^n\) in two ways: as a product of one-variable integrals, and in polar coordinates.
\item[\ref{pr:base-change}] A basis of \(V\) over \(K\) is a basis of \(V \otimes_K K'\).  Use \(\sum_{w \mid v} [K'_w : K_v] = [K':K]\).
\end{description}
\paragraph{Section~\ref{sec:HS}.}
\begin{description}
\item[\ref{pr:HS-P1-check}] Here \(d = 1\).  For the second statement use Problem~\ref{sm-pr:deg-basic}.
\item[\ref{pr:L2-sup}] Diagonalize one norm with respect to the other.  Two smooth volume forms on a compact manifold are bounded by constant multiples of each other.
\end{description}
\paragraph{Section~\ref{sec:definite}.}
\begin{description}
\item[\ref{pr:def-deg}] The forms make up an orthogonal sum over \(\Cl_n\).  Scaling a norm by \(\lvert\Gamma_L\rvert^{-1/2}\) changes the degree as in Problem~\ref{sm-pr:deg-basic}, with \(n_v = 2\).
\item[\ref{pr:def-n1}] \(\Gamma_{I_1} = \mu_4\), and \(\lVert e^a \rVert = 1\).
\item[\ref{pr:sym-sphere}] Problem~\ref{pr:PN}.
\item[\ref{pr:def-n2}] The diagonal subgroup acts on \(e^\nu\) by \(\ii^{\nu_1 j_1 + \nu_2 j_2}\).  Then use Problem~\ref{pr:def-deg}.
\item[\ref{pr:def-n2-asym}] The rank is \(\lfloor a/8 \rfloor + 1\).  The sum of \(\log\binom{a}{\nu_1}\) over \(\nu_1 \equiv 0 \pmod 4\), \(\nu_1 \le a/2\), is about one eighth of the full sum, which is computed in Problem~\ref{pr:degED-asym}; since \(k \mapsto \log\binom{a}{k}\) increases on \([0, a/2]\), the error is at most \(\log\binom{a}{a/2} \le a\log 2\).
\end{description}
\paragraph{Section~\ref{sec:uniform}.}
\begin{description}
\item[\ref{pr:rank-ratio}] Write the ratio as a product of \(n - 1\) factors.
\item[\ref{pr:trivial-lower}] The monomial basis and Problem~\ref{pr:PN}.
\item[\ref{pr:level}] Problem~\ref{sm-pr:deg-basic}.
\end{description}

\begin{thebibliography}{MB1}
\bibitem[AB]{AB95} A.~Abbes and T.~Bouche, \emph{Th\'eor\`eme de Hilbert--Samuel ``arithm\'etique''}, Ann.\ Inst.\ Fourier \textbf{45} (1995), 375--401.
\bibitem[Adr]{Adr} J.~Adriazola, \emph{The Alp\"oge proof of the Riemann hypothesis}, manuscript, revised 24 November 2026, \url{https://github.com/jimjamzola/RH-Alpoge-Reconstruction}.
\bibitem[Alp]{Alp} L.~Alp\"oge, untitled note announcing a proof of the Riemann hypothesis, 20 November 2026, \url{https://alpo.ge/remark.pdf}.
\bibitem[BGS]{BGS} J.-B.~Bost, H.~Gillet and C.~Soul\'e, \emph{Heights of projective varieties and positive Green forms}, J.\ Amer.\ Math.\ Soc.\ \textbf{7} (1994), 903--1027.
\bibitem[DPh]{DPh} S.~David and P.~Philippon, \emph{Minorations des hauteurs normalis\'ees des sous-vari\'et\'es des tores}, Ann.\ Scuola Norm.\ Sup.\ Pisa Cl.\ Sci.\ (4) \textbf{28} (1999), 489--543.
\bibitem[GS1]{GS90} H.~Gillet and C.~Soul\'e, \emph{Arithmetic intersection theory}, Publ.\ Math.\ IH\'ES \textbf{72} (1990), 93--174.
\bibitem[GS2]{GS92} H.~Gillet and C.~Soul\'e, \emph{An arithmetic Riemann--Roch theorem}, Invent.\ Math.\ \textbf{110} (1992), 473--543.
\bibitem[Zh]{Zh} R.~Zhang, \emph{A preliminary note on Alp\"oge's proposed proof of the generalized Riemann hypothesis}, preprint, 2026, hexagon:2611.00809.
\end{thebibliography}

\end{document}
